\chapter{Market Participants and Their Incentives}
\label{ch:participants}
\chaptermark{Market participants}

Every trade has two sides. When your buy fills, \emph{someone} sold;
when your sell fills, \emph{someone} bought. That someone's motivation,
information, and pressures determine whether your trade was a good
idea. The same resting quote is a free gift if the counterparty is an
inattentive retail trader dumping a position, and a lethal trap if it
is an insider with next-quarter's earnings release. Until you can
reason about \emph{who is on the other side}, you cannot reason about
your own edge.

This chapter introduces the seven categories of participant you will
meet repeatedly: why they trade, how their orders appear in the limit
order book, and how to recognise them in market data. The centrepiece
is the distinction between \textbf{informed} and \textbf{uninformed}
traders --- the foundation of every microstructure model in
\cref{ch:spreads_roll,ch:kyle_gm}, and the reason a market maker can
earn money posting quotes, or lose it quickly if they misjudge who is
hitting them.

\begin{backgroundbox}[Prerequisites]
From \cref{ch:lob}:
\begin{itemize}
  \item The \textbf{limit order book} as two sorted lists of resting
  buy/sell orders.
  \item \textbf{Limit} (post a price) vs \textbf{market} (cross the
  book) orders; the \textbf{order lifecycle}: submit, queue, match,
  report.
  \item \textbf{Makers} (resting order) and \textbf{takers} (marketable
  order crossing). Exchanges charge the taker a fee and pay the maker
  a rebate.
  \item \textbf{Bid}, \textbf{ask}, \textbf{spread}, \textbf{mid},
  \textbf{tick size}.
\end{itemize}
No prior finance knowledge assumed beyond these mechanics.
\end{backgroundbox}

\section{The trader taxonomy}
\label{sec:taxonomy}

Larry Harris's \emph{Trading and Exchanges} \citeHarris{} organises
market participants by why they trade, not what they trade. Motivation
predicts behaviour: a participant trading on information acts
differently from one rebalancing a pension fund, even sending identical
orders to the same exchange. The seven categories below are not
mutually exclusive --- one firm can run desks in several roles --- but
they map cleanly onto the strategies in
\cref{ch:market_making,ch:arbitrage_zoo}.

\subsection{Market makers}
\label{subsec:mm}

\begin{definition}[Market maker]
\index{market maker|textbf}
A \emph{market maker} is a participant whose strategy is to
continuously quote both a bid and an ask on a product, profiting from
the spread $\spreads = \askp - \bidp$ between the two by being on the
maker side of many trades in both directions. A pure market maker does
not take a view on the direction of the underlying price; their goal
is to end the trading session with roughly zero net position
(\emph{flat}) and cash equal to their aggregated spread capture minus
losses.
\end{definition}

The canonical mental picture is the airport currency booth: buy
dollars from tourists at one rate, sell at a higher rate, and if flow
balances, bank the difference. Electronically the market maker posts a
buy limit at $\bidp$ and a sell limit at $\askp$, then waits. When a
taker hits either side, position moves one unit and cash moves by the
fill price. If a buy and a sell fill in close succession, they bought
at $\bidp$ and sold at $\askp$, netting $\spreads$ per unit at zero
end-of-session exposure.

The behaviour is recognisable in the LOB. A market maker refreshes
quotes on both sides, clustered near the top of the book in
competition with other market makers, cancelling and re-posting as the
mid moves. They rarely send a \emph{marketable} order (one that
crosses the book) except to unwind a lopsided inventory.
\Cref{ch:market_making} derives the optimal quoting strategy under
inventory risk; for now, the takeaway is that market makers are
maker-side, two-sided, and almost always present.

\paragraph{Identification heuristic.} Very high fill count, near-equal
split between buy and sell fills, tiny average holding time, positions
that oscillate near zero, quotes that appear within one tick of the
best on both sides.

\subsection{Arbitrageurs}
\label{subsec:arb}

\begin{definition}[Arbitrageur]
\index{arbitrageur|textbf}
An \emph{arbitrageur} is a participant who trades simultaneously in
two or more markets (or two or more products) to lock in a price
discrepancy. A pure \emph{riskless} arbitrage requires the discrepancy
to close with certainty; a \emph{statistical} arbitrage merely
requires it to close in expectation.
\end{definition}

The archetype is triangular FX arbitrage: if EUR/USD/JPY cross-rates
imply a loop EUR~$\to$~USD~$\to$~JPY~$\to$~EUR worth more than one
euro, an arbitrageur sends three market orders in a race against
everyone else doing the same. Each leg takes liquidity from a
different book; the profit is the loop's excess; the risk is that a
leg fails to fill at the expected price, leaving the position
directionally exposed.

In the LOB, arbitrageurs are almost exclusively \textbf{takers}: the
opportunity closes in milliseconds, too fast to rest a limit order.
Their flow is \emph{correlated across venues} --- simultaneous hits on
EUR/USD at two venues are likely one arbitrageur firing two legs.
\Cref{ch:arbitrage_zoo} catalogues the arbitrage types (triangular,
cash-and-carry, put--call parity, ETF creation--redemption, and more);
for now, note that arbitrageurs force prices across venues to stay in
line.

\paragraph{Identification heuristic.} Short bursts of taker orders,
highly correlated timestamps across multiple venues or products,
near-zero inventory held over the round, directionless aggregate P\&L
(they do not care which way the market moves).

\subsection{Informed traders}
\label{subsec:informed}

\begin{definition}[Informed trader]
\index{informed trader|textbf}
An \emph{informed trader} is a participant who possesses private
information about the fair value of a product that is not yet
reflected in the market price. They trade on that information: buying
when their private signal says the product is underpriced, selling
when it says overpriced, and exiting the position once the price moves
to their target.
\end{definition}

The prototypical informed trader in academic models is an
\emph{insider}: an executive who knows an earnings announcement is
coming, an analyst who has done original research, a trader who has
seen a regulatory filing minutes before it goes public. Harris's
broader sense of ``informed'' covers any participant whose expected
value for the product exceeds the market's information set --- not
only illegal insider trading but also superior models, proprietary
data feeds, or faster news parsing.

From the market maker's point of view, they \emph{cannot see the
information nor the label on the order}. An informed buy and a noise
buy are indistinguishable at the message level; only the subsequent
price move reveals which hit them. This is the setup for both
classical theories of the bid--ask spread in \cref{ch:kyle_gm}:
Glosten--Milgrom models informed trading as an adverse-selection tax
on the market maker, and Kyle's 1985 model derives the equilibrium in
which a single strategic informed trader hides inside noise flow while
the market maker learns from order flow alone.

\paragraph{Identification heuristic.} Consistently on the same side of
a subsequent price move; buys precede up-moves and sells precede
down-moves with far higher hit-rate than chance. The hardest
counterparty to recognise in real time; the easiest in hindsight.

\subsection{Noise traders (liquidity traders)}
\label{subsec:noise}

\begin{definition}[Noise trader]
\index{noise trader|textbf}
A \emph{noise trader}, also called a \emph{liquidity trader} or
\emph{uninformed trader}, is a participant whose reason for trading is
unrelated to private information about fair value. They trade because
they need cash, because a pension fund is rebalancing on a calendar
schedule, because a mutual fund is meeting redemptions, because
someone pressed the wrong button, or because an automated strategy's
signal was pure statistical nonsense.
\end{definition}

The terminology is historical and slightly unfair: a pension fund
rebalancing a multi-billion-dollar portfolio is not literally
\emph{noisy}, just orthogonal to information about the next price
move. What matters is that the noise trader's decision is
\emph{uncorrelated with the product's true value}: on average they
buy as often before a down-move as before an up-move. To the market
maker, the noise trader is the perfect counterparty --- every fill is
a fair bet and the spread collected is profit.

Real markets always contain a mix of informed and noise flow, and
every microstructure model in this book is built around that mix. The
market maker's willingness to quote a tight spread depends on the
\emph{noise fraction} being high; when informed flow dominates,
spreads widen or market makers withdraw. \Cref{ch:spreads_roll} makes
this quantitative.

\paragraph{Identification heuristic.} Uncorrelated with subsequent
price moves; roughly 50\% of trades face up-moves and 50\% face
down-moves. Often larger order sizes and longer holding horizons than
information-driven flow. In Prosperity, the scripted ``retail'' bots
that simply hit the top of the book at random intervals are pure
noise traders.

\subsection{High-frequency traders (HFTs)}
\label{subsec:hft}

\begin{definition}[High-frequency trader]
\index{high-frequency trader (HFT)|textbf}\index{HFT|textbf}
A \emph{high-frequency trader (HFT)} is a participant whose strategy
operates at microsecond to sub-millisecond latencies, rewriting quotes
hundreds or thousands of times per second in response to every incoming
message on the book. HFTs are technologically, not fundamentally,
defined: the category includes market-making HFTs, arbitrage HFTs, and
short-horizon signal-driven HFTs.
\end{definition}

HFTs sit on top of the taxonomy rather than inside it: most are a
hybrid of market-making and arbitrage roles, operating at a latency
regime only hardware-colocated strategies can reach. Their defining
feature is intense concern with \textbf{queue position}: in a
price--time priority market (\cref{ch:lob}) the first order at a price
level fills first, and moving from tenth- to first-in-queue can double
a strategy's fill rate. HFTs pay to co-locate servers in the same data
centre as the matching engine, measured in metres of fibre, so their
messages arrive ahead of competitors'.

The economic contribution of HFTs is contested. The defensive case:
HFT market making has driven spreads on liquid products to historic
lows, benefiting every other participant. The critical case: HFTs
impose an adverse-selection tax on slower participants by racing in
front of institutional flow. \Cref{ch:order_flow} develops the
order-flow imbalance framework underlying most short-horizon HFT
signals.

\paragraph{Identification heuristic.} Enormous message rates (cancels
per second), near-zero average holding time, aggressive quote
competition at the top of the book, sensitivity to microsecond-level
latency events. In exchange-level data they show up as the firms
submitting 95\%+ of all cancellations.

\subsection{Institutional asset managers}
\label{subsec:inst}

\begin{definition}[Institutional asset manager]
\index{institutional asset manager|textbf}
An \emph{institutional asset manager} is a firm that trades on behalf
of long-horizon capital: pension funds, mutual funds, sovereign wealth
funds, university endowments, insurance companies, and the like. Their
positions are measured in hundreds of millions to hundreds of billions
of dollars, their holding periods run from months to decades, and
their trading is driven by portfolio construction mandates rather than
by microstructure signals.
\end{definition}

Institutional flow is the natural counterparty of the HFT market
makers above: slow, size-constrained orders meet fast, two-sided
quotes. The institutional manager consumes liquidity (takes); the
HFT provides it (makes), with the maker--taker fee schedule
(\cref{ch:lob}) aligning the incentives.

An institutional manager buying 10 million shares cannot send one
market order for 10 million: it would eat the entire visible book and
move price by percent. Instead they slice the parent order over hours
or days using execution algorithms such as VWAP, TWAP, and
participation-of-volume targets. \Cref{ch:almgren_chriss} derives the
optimal schedule under a linear impact model. For now: institutional
trading is \emph{size-constrained and execution-sensitive} --- P\&L
is dominated by the price received on each slice, not by the raw
direction of the trade.

To the market maker, institutional flow is \emph{partially informed}:
the rebalance decision is uncorrelated with short-horizon moves
(looks like noise), but the signed persistence of a parent sliced
over hours reveals its direction (looks like information). This is
why institutional orders leak impact even when executed carefully ---
persistent one-sided flow is itself informative.

\paragraph{Identification heuristic.} Persistent one-sided flow over
hours with roughly constant slice size; use of VWAP/TWAP child orders
that match the day's volume profile; the parent order ends when the
target quantity is reached rather than when a price target is hit.

\subsection{Retail traders}
\label{subsec:retail}

\begin{definition}[Retail trader]
\index{retail trader|textbf}
A \emph{retail trader} is an individual trading their own personal
account through a broker. Order sizes are small (typically tens to
thousands of dollars per trade), the trader is usually less informed
than professional participants, and their flow reaches the exchange
only after the broker has first had an opportunity to internalise it
or to sell it to a wholesaler.
\end{definition}

Retail motivations are varied and typically non-professional:
long-horizon saving for retirement, speculation on a view picked up
from news or social media, hobbyist day-trading, or hedging a
concentrated stock grant from an employer. What unites the category
is that trading is not the participant's job and the capital at risk
is personal.

The retail trader never reaches the exchange directly. Their order
travels first to their broker, who in most modern markets either
executes against the broker's own inventory (\emph{internalisation})
or sells the order to a wholesale market maker for a small fee under
\emph{payment for order flow (PFOF)}. The wholesaler then matches it
internally or routes to the exchange. By the time a retail order
hits the public book, it has been inspected by at least one
professional.

Wholesale market makers pay for retail flow because it is, on average,
heavily noise-dominated. A retail trader is more likely selling to
pay a bill than because they have private information about earnings.
Paying for retail flow therefore buys the best counterparty in the
market --- the noise trader of \cref{subsec:noise}. Retail flow is
not uniformly uninformed: retail-level insider trading, meme-stock
coordination, and reactions to public news all inject short bursts of
information into the channel.

\paragraph{Identification heuristic.} Small average order size,
heavy concentration in a handful of popular tickers, bursts of
correlated buying or selling around news events, and, in venues where
it is observable, the characteristic latency signature of a broker's
customer router rather than a direct exchange feed.

\section{Informed vs uninformed: the single most important distinction}
\label{sec:informed_vs_uninformed}

The seven categories collapse onto a single axis that dominates all
others: is the counterparty \emph{informed} or \emph{uninformed}? A
market maker trades with noise all day and flees informed flow; an
execution algorithm disguises its parent order as noise; an HFT
arbitrageur wants to act on information first. Every microstructure
result in \cref{ch:spreads_roll,ch:kyle_gm} is ultimately a statement
about how this distinction shapes prices.

\begin{definition}[Informed and uninformed flow]
\index{informed flow|textbf}\index{uninformed flow|textbf}
Let $V$ denote the unobserved fair value of a product a short time
after the current instant. An order is \emph{informed} if the
participant submitting it has information that changes their
conditional expectation of $V$ relative to the public information set.
An order is \emph{uninformed} (equivalently \emph{noise} or
\emph{liquidity}) if its submission is conditionally independent of
$V$ given public information: the participant is trading for reasons
orthogonal to the product's short-horizon value.
\end{definition}

The market maker \emph{cannot see this label}. From the message
stream, an informed buy and a noise buy are bit-for-bit identical:
side, quantity, timestamp, maybe an anonymised account identifier.
The only way to infer the label is to observe what the price does
\emph{after} the trade. Over many trades the market maker can build a
statistical model of each counterparty's informativeness; on a single
trade, they are flying blind.

This blindness is asymmetric against the market maker. Fills against
noise traders are fair coin flips around the mid, and the spread
turns them into profit. Fills against informed traders are
systematically adverse: the informed buyer only hits the ask when
fair value strictly exceeds the ask, leaving the market maker's
post-trade inventory underwater. This is \textbf{adverse
selection}.

\subsection{A toy adverse-selection calculation}
\label{subsec:toy_adverse}

Consider the simplest model. A product has two possible fair values,
$V = 100$ and $V = 110$, equally likely a priori. Fraction $\alpha =
0.3$ of arriving traders are \emph{informed}: they know $V$ and buy
when $V = 110$, sell when $V = 100$. The other $1 - \alpha = 0.7$ are
\emph{noise}: they buy or sell with probability $\tfrac{1}{2}$ each,
independent of $V$. The market maker posts a single price at the mid,
$\midp = 105$, and trades with whoever arrives.

Suppose a buy order arrives and fills at $105$. What does the market
maker expect to lose?

\paragraph{Step 1: $P(\text{buy} \mid V)$.}
\[
  P(\text{buy} \mid V = 110) \;=\; \underbrace{\alpha \cdot 1}_{\text{informed, buys}}
    \;+\; \underbrace{(1 - \alpha) \cdot \tfrac{1}{2}}_{\text{noise, random}}
    \;=\; 0.3 + 0.7 \cdot 0.5 \;=\; 0.65.
\]
\[
  P(\text{buy} \mid V = 100) \;=\; \underbrace{\alpha \cdot 0}_{\text{informed, sells}}
    \;+\; \underbrace{(1 - \alpha) \cdot \tfrac{1}{2}}_{\text{noise, random}}
    \;=\; 0 + 0.35 \;=\; 0.35.
\]

\paragraph{Step 2: $P(\text{buy})$.}
\begin{align*}
  P(\text{buy}) &\;=\; P(V = 110)\,P(\text{buy} \mid V = 110)
                     + P(V = 100)\,P(\text{buy} \mid V = 100) \\
                &\;=\; 0.5 \cdot 0.65 + 0.5 \cdot 0.35
                 \;=\; 0.50.
\end{align*}

\paragraph{Step 3: posterior over $V$ given a buy, by Bayes.}
\[
  P(V = 110 \mid \text{buy})
    \;=\; \frac{P(V = 110)\,P(\text{buy} \mid V = 110)}{P(\text{buy})}
    \;=\; \frac{0.5 \cdot 0.65}{0.50}
    \;=\; 0.65.
\]
\[
  P(V = 100 \mid \text{buy}) \;=\; 1 - 0.65 \;=\; 0.35.
\]

\paragraph{Step 4: expected fair value given a buy.}
\[
  \E[V \mid \text{buy}]
    \;=\; 0.65 \cdot 110 + 0.35 \cdot 100
    \;=\; 71.5 + 35 \;=\; 106.5.
\]

\paragraph{Step 5: market maker's expected P\&L on this fill.}
The market maker sold at $\midp = 105$ and now holds $-1$ unit of a
product whose conditional fair value is $106.5$. The expected loss
per unit is
\[
  \underbrace{105}_{\text{price received}}
    \;-\; \underbrace{106.5}_{\E[V \mid \text{buy}]}
    \;=\; -1.5.
\]
The market maker loses $1.5$ on average on every buy that hits their
ask, because the buy is informationally tilted toward the high-value
state. By symmetry, every sell that hits their bid loses $1.5$ as
well: sells come preferentially from informed traders who know $V =
100$, so the market maker buys at $105$ a product worth $103.5$,
giving P\&L $= -1.5$ per unit.

The market maker cannot stay at the mid. Allowing a genuine spread
--- ask strictly above $105$, bid strictly below --- shrinks losses,
and at a wide enough spread turns them positive. ``How wide must the
spread be to offset adverse selection?'' is the Glosten--Milgrom
problem, derived in full in \cref{ch:kyle_gm}. The toy calculation
here shows that \emph{a nonzero spread is forced by the mere
existence of informed flow}, even without order-processing costs or
inventory risk. The other two spread drivers are added in
\cref{ch:spreads_roll}.

\begin{keyfactbox}[The market maker's fundamental problem]
The market maker sees only anonymous order flow and cannot tell
informed from noise. Informed flow is systematically adverse (buyers
buy only when fair value exceeds the ask, sellers sell only when fair
value is below the bid), so fills against it lose money in
expectation. The whole point of quoting a spread rather than a single
price is to charge noise flow enough to cover losses to informed
flow. This tension is \textbf{adverse selection}, the first-order
phenomenon every microstructure model quantifies.
\end{keyfactbox}

\begin{intuitionbox}
You would never play poker against someone who can see your hand. A
market maker plays a worse game: against a table of strangers, some
of whom see the hand and some of whom do not, and the market maker is
not told which. Their only tool is to charge enough per hand to break
even averaged over the whole table. That charge is the spread.
\end{intuitionbox}

\begin{pitfallbox}[``Small spread = cheap trading cost'' is wrong for makers]
Novices assume a tight spread means cheap trading. True for a
\emph{taker} (\cref{ch:lob}): tight spread means small crossing
distance. False for a \emph{maker}: tight spread means small reward
per fill while adverse-selection losses continue unchanged. On a
product with heavy informed flow, a tight spread is a warning, not an
invitation. \Cref{ch:market_making} shows how a market maker skews
quotes \emph{inside} or \emph{outside} the top of book depending on
their view of adverse selection.
\end{pitfallbox}

\section{Maker/taker economics in full}
\label{sec:maker_taker_econ}

\Cref{ch:lob} introduced the maker--taker fee schedule mechanically:
the exchange charges the taker a small fee, pays the maker a small
rebate, and keeps the difference. This chapter can now explain
\emph{why}, and why the asymmetry is necessary.

The maker of every trade is someone whose limit order was resting when
the trade happened. That order has been exposed to \textbf{adverse
selection} for every microsecond it sat: any informed trader arriving
during that window would have picked off the stale quote. The longer a
quote rests, the more likely the next fill is adverse. Posting is
costly, and someone must compensate the posters or they will stop
posting.

The taker, in the same trade, crossed the book. They paid the spread
--- the gap between fill price and mid --- and received the one thing
the market maker cannot promise: \textbf{immediacy}. The taker traded
right now, at the quoted price, without risk of the opportunity
disappearing. This is valuable to any participant with a
short-shelf-life trade: an arbitrageur whose loop closes in
milliseconds, an execution algorithm needing to complete by the close,
a retail trader reacting to a news event. Takers pay for immediacy
because without it their strategy does not exist.

\begin{keyfactbox}[The economic symmetry of the maker/taker fee]
The maker bears adverse-selection risk, inventory risk, and the
opportunity cost of capital in resting orders. In return they get (a)
the spread and (b) an exchange rebate. The taker gets immediacy. In
return they pay (a) the spread and (b) an exchange fee. The exchange
clips the difference to fund matching, compliance, and margin.
Without the rebate, no posted liquidity; without the fee, no
exchange.
\end{keyfactbox}

A market charging neither side would still function in principle ---
participants would trade for the spread alone --- but the exchange
would have no revenue, and without a rebate the incentive to post
rather than wait-and-take disappears. Rebates are typically a few
tenths of a basis point, but at the scale of rebate-harvesting market
making (\cref{ch:market_making}) they are the entire business model.

Not every venue uses the same schedule. A \textbf{maker--maker model}
rebates both sides and funds the exchange from subscription and data
fees. Some derivatives exchanges use \textbf{pro-rata allocation}
instead of price--time priority, devaluing queue position and
shifting market-making economics. Dark pools charge neither side and
exist as venues for institutional block crosses that avoid signalling.
\Cref{ch:market_making} treats the full range.

\section{Dealer markets vs exchange markets}
\label{sec:dealer_vs_exchange}

So far this book has described an \textbf{order-driven market}: a
central limit order book with anonymous participants, price--time
priority, and a matching engine that pairs crossing orders. This is
the dominant structure for listed equities, listed futures, major
cryptocurrencies, and all of Prosperity. It is not the only
architecture, and several of the world's largest asset classes do
\emph{not} trade this way.

\begin{definition}[Order-driven market]
\index{order-driven market|textbf}
An \emph{order-driven market} is one in which all participants submit
limit and market orders to a single shared limit order book, trades
are matched by an anonymous matching engine under fixed priority
rules, and every participant faces the same book. Prices emerge from
the interaction of orders; no single participant is privileged.
Examples: equities on NASDAQ and NYSE, futures on CME, cryptocurrencies
on Binance, and Prosperity.
\end{definition}

\begin{definition}[Quote-driven or dealer market]
\index{dealer market|textbf}\index{quote-driven market|textbf}
A \emph{quote-driven market} (also called a \emph{dealer market}) is
one in which a participant wishing to trade does not send an order to
a central book but instead sends a \emph{request for quote (RFQ)} to
one or more dealers, typically large banks. Each dealer responds with
a bid and an ask at which they are willing to trade; the client picks
the best one and trades against that dealer's inventory. The dealer
takes the position onto their own book and is responsible for
hedging it. There is no anonymous matching; the dealer is a
\emph{counterparty}, not a neutral intermediary. Examples: most over-
the-counter (OTC) bond markets, interest rate swaps, FX spot for
institutional size, OTC derivatives.
\end{definition}

The difference is fundamental enough that this book's intuitions
transfer only partially. In an order-driven market the sell-side bank
is just another participant; in a quote-driven market the bank is the
\emph{only} participant on one side, with entirely different
incentives. A dealer pricing an interest rate swap is not a neutral
matching engine: they price against their own inventory, their fair
value, their hedging costs, and their assessment of how informed the
client is. A macro hedge fund gets quoted wider; a corporate treasurer
rolling a hedge gets quoted tighter.

The microstructure models of \cref{ch:spreads_roll,ch:kyle_gm} were
developed for order-driven markets, but the core intuition ---
adverse-selection pressure on the liquidity provider --- carries over.
A dealer faces the same problem as a market maker with a different
interface: they see the client's identity (helps: classify the flow)
but commit to a price before the client commits to trade (hurts: the
client shops around). \Cref{ch:strats_role} returns to the day-to-day
life of a Strat on the quote-driven side.

\begin{gsbox}
On a Goldman rates desk, a pension fund calls asking a price on a
\$500m 10-year interest rate swap. There is no central book for swaps
in that size. You --- the Strat --- receive the RFQ; your system
computes a fair mid from the current swap curve; you add a spread
reflecting your inventory (already long duration? charge more), your
view of the client's informedness (routine corporate hedge or macro
fund taking a directional view?), and your hedging cost in the linear
futures market where you will offset delta within minutes. The client
picks your price or someone else's; if yours, the position lands on
your book and you immediately sell Eurodollar futures to neutralise
first-order rate exposure. \Cref{ch:strats_role} teaches how every
piece of this workflow is implemented.
\end{gsbox}

The boundary is not watertight. Some equity venues run both: a
central book by day and an RFQ facility for block trades. Some crypto
products trade on central books and via OTC dealers. Some order-driven
exchanges host dealer-style market makers (``lead market makers'' in
ETFs, designated market makers on NYSE, specialists in older designs).
What does not change is the economic pressure: the liquidity provider
--- central-book market maker or sell-side dealer --- is paid to
absorb risk from flow they cannot fully classify. The game is the
adverse-selection game of \cref{ch:kyle_gm}.

\section{The Olivia story: an informed trader in the wild}
\label{sec:olivia}

The informed/uninformed distinction can feel abstract. Prosperity~3
provided a vivid real example showing the informed-trader category is
a recognisable pattern in actual order flow --- and exploiting it
paid.

\begin{prosperitybox}
In the 2025 IMC Prosperity~3 competition, participants traded against
scripted bots whose trades were \emph{anonymous} for the first four
rounds: you saw trade size, price, and direction but not which bot
acted. In Round~5 IMC made historical \textbf{trader IDs} public,
lifting the anonymity curtain and letting teams tag every past trade
with its submitting bot.

One bot, \textbf{Olivia}, had followed a strikingly simple rule
throughout: on Squid Ink, buy $15$ lots at the daily low, sell $15$
lots at the daily high. In the vocabulary of
\cref{sec:informed_vs_uninformed} she was an archetypal informed
trader, acting on private information (the day's price extrema) no
one else could see until after the fact.

The Frankfurt Hedgehogs \citep{frankfurtHedgehogs2025imcP3}, who
finished 2nd overall, spotted this during the first four rounds
without knowing the bot's name: a particular anonymous counterparty
consistently bought at what turned out to be the daily minimum and
sold at the daily maximum, a pattern their running-extremum detector
could latch onto. The optimal strategy on Squid Ink became ``mimic the
informed trader'': buy when Olivia buys, sell when Olivia sells. This
single signal generated roughly $8{,}000$ SeaShells on Squid Ink in
Round~1. In Round~2 they used Olivia's inferred Croissants position
(a picnic-basket constituent) to tilt basket-spread thresholds ---
the cross-product extension of \cref{ch:basket_etf_arb}.

When Round~5 revealed trader IDs, the Hedgehogs already had a working
detector; the reveal merely let them swap heuristic matching for a
literal trader-ID check, eliminating false positives. Several other
top-ten teams converged on the same strategy.

The story maps onto this chapter's theory. Olivia was informed in the
strict sense: a private signal, acted on. Any team running pure
market-making on Squid Ink was unwittingly the market maker of
\cref{sec:informed_vs_uninformed}, paying for adversely-selected flow.
Teams that identified and copied Olivia were on the informed side, and
P\&L flowed accordingly. Lesson for Prosperity~4 (and for any real
exchange with mixed counterparties): the informed/uninformed
distinction is not abstraction. It is the difference between making
and losing money, visible in the order flow if you know where to
look.
\end{prosperitybox}

\section{Key facts to memorise}
\label{sec:ch2_key_facts}

\begin{itemize}
  \item The seven participant categories --- \textbf{market makers},
  \textbf{arbitrageurs}, \textbf{informed traders}, \textbf{noise
  (liquidity) traders}, \textbf{HFTs}, \textbf{institutional asset
  managers}, \textbf{retail traders} --- are distinguished by
  motivation, not by the orders they send.

  \item The most consequential split is \textbf{informed vs
  uninformed}: flow correlated with subsequent price moves vs flow
  orthogonal to them. Every result in
  \cref{ch:spreads_roll,ch:kyle_gm} rests on this distinction.

  \item \textbf{Adverse selection}: a market maker's fills against
  informed counterparties systematically lose money because the
  informed trader crosses only when fair value is favourable. In the
  toy model of \cref{subsec:toy_adverse} ($\alpha = 0.3$, $V \in
  \{100, 110\}$, mid $105$), expected loss per fill is $\$1.5$ per
  unit, forcing a nonzero spread.

  \item The \textbf{maker--taker fee schedule} is economic symmetry:
  makers get a rebate for bearing adverse selection and inventory
  risk; takers pay a fee for immediacy. Without the rebate, no one
  posts; without the fee, the exchange has no revenue.

  \item \textbf{Order-driven} (CLOB, anonymous, price--time priority)
  and \textbf{quote-driven / dealer} (RFQ, dealer holds inventory)
  markets are architecturally distinct. Equities, futures, crypto,
  Prosperity sit in the first; bonds, swaps, institutional FX spot,
  OTC derivatives sit in the second. The adverse-selection game is
  the same.

  \item \textbf{Queue position} matters under price--time priority;
  this is the economic reason HFTs co-locate.

  \item \textbf{Institutional flow} is partially informed: slices are
  uncorrelated with the next tick, but parent-order persistence leaks
  direction. \Cref{ch:almgren_chriss} covers large-order execution.

  \item \textbf{Retail flow is valuable because it is uninformed}:
  wholesalers pay brokers for PFOF access, because a noisy counterparty
  is the best counterparty.

  \item The \textbf{Olivia story} is empirical ground truth for the
  informed-trader category --- detectable from order flow alone, and
  top-ten teams copied her for direct P\&L. The informed/uninformed
  distinction is a pattern in real data, not an abstraction.
\end{itemize}
